\section{Conclusions}
\label{sec:conclusion}

We presented MgfrOFR, a Mamba-based old film restoration framework that augments MambaOFR with explicit multi-granularity feature balancing. The Multi-Granularity Feature Extractor (MGFE) reorganizes propagated features into coarse-to-fine restoration-token branches and recalibrates the spatial stream through channel-wise feedback, while Ranking-Aware Triplet Regularization (RATR) discourages branch collapse by penalizing correlated temporal-similarity rankings during training.

With the final average-fusion configuration on the synthetic benchmark, MgfrOFR achieves the strongest full-reference scores among the tested baselines, improving PSNR by 1.02~dB over MambaOFR (27.86~dB vs.\ 26.84~dB). On 216 real-world archival clips, MgfrOFR obtains the lowest NIQE (5.12), improves BRISQUE over MambaOFR (24.79 vs.\ 29.79), and shows a lower CLIP-IQA+ value (0.4092) than the sharper enhancement-oriented references and old film restoration baselines. Because the real-world setting lacks ground truth, these no-reference gains should be interpreted as evidence of improved statistical naturalness rather than definitive proof of perceptual superiority. The ablation study further shows that MGFE and RATR are complementary: MGFE alone improves the baseline but remains weaker than the full model on full-reference perceptual metrics, whereas the full model gives the best combined result when branch diversity is explicitly enforced through RATR and average fusion.

These results suggest that representation-level branch balance is a useful inductive bias for heterogeneous-degradation video restoration and practical archival digitization workflows. The main implication is that feature-space symmetry and temporally coordinated ranking behavior can improve restoration under degradation-induced asymmetry without changing the baseline inference interface. Several limitations remain: the method inherits the synthetic-to-real domain gap of REDS-based training; RATR adds a training-time cost that grows approximately quadratically with the batch-time descriptor set, increasing training time by approximately 7.6\% over MambaOFR under our setting; and evaluation is limited to the SRWOV benchmark and no-reference proxies on authentic footage. Future work will focus on adaptive branch activation conditioned on degradation severity, frequency-aware branch specialization for structural recovery and texture restoration, lighter ranking constructions such as selective anchor sampling, human-in-the-loop evaluation, and broader validation on additional archival benchmarks.
